\documentclass[../../main.tex]{subfiles}
\begin{document}

{\bf [解答]}

\begin{figure}[htb]
  \centering
  \begin{tikzpicture}[x=1.8cm, y=0.75cm, >=stealth,
    every node/.style={font=\small}]

    % --- 1. 定数・座標 ---
    % b = 2.5, a = -0.8125
    \def\a{-0.8125}
    \def\b{2.5}
    \pgfmathsetmacro{\xV}{-\b / (2 * \a)}       % 1.538
    \pgfmathsetmacro{\yV}{-\b * \b / (4 * \a)}  % 1.923
    \pgfmathsetmacro{\xR}{-\b / \a}             % 3.077

    \coordinate (O) at (0, 0);
    \coordinate (V) at (\xV, \yV);
    \coordinate (R) at (\xR, 0);
    \coordinate (LineX) at (2.5, 0);
    \coordinate (LineY) at (0, 5);

    % --- 2. 面積 S の塗り ---
    \fill[blue!12, domain=0:\xR, samples=60]
    plot (\x, {\a*\x*\x + \b*\x}) -- (O) -- cycle;
    \node[blue!85!black, font=\normalsize] at (0.5*\xV, 0.45*\yV) {$S$};

    % --- 3. 座標軸 ---
    \draw[->, thick, gray!80!black] (-0.4, 0) -- (3.6, 0) node[right] {$x$};
    \draw[->, thick, gray!80!black] (0, -0.6) -- (0, 5.5) node[above] {$y$};
    \node[below left=2pt] at (O) {$O$};

    % --- 4. 直線 y = -2x + 5 ---
    \draw[line width=1.1pt, teal!75!black, domain=-0.2:2.8, samples=2]
    plot (\x, {-2*\x + 5});
    \node[teal!85!black, font=\footnotesize, above right] at (0.2, 4.6) {$y = -2x + 5$};

    % --- 5. 放物線 y = ax^2 + bx ---
    \draw[line width=1.4pt, blue!85!black, domain=-0.2:3.3, samples=70]
    plot (\x, {\a*\x*\x + \b*\x});
    \node[blue!85!black, font=\footnotesize, above right] at (2.2, 1.8) {$y = ax^2 + bx$};

    % --- 6. 頂点 V の投影線とラベル ---
    \draw[dashed, gray!60] (\xV, 0) -- (V) -- (0, \yV);
    \node[below=2pt, font=\footnotesize] at (\xV, 0) {$-\frac{b}{2a}$};
    \node[left=2pt, font=\footnotesize] at (0, \yV) {$-\frac{b^2}{4a}$};

    % x切片
    \node[below=2pt, font=\footnotesize] at (\xR, 0) {$-\frac{b}{a}$};

    % 直線の切片
    \node[left=2pt, font=\footnotesize] at (LineY) {$5$};
    \node[below=2pt, font=\footnotesize] at (LineX) {$\frac{5}{2}$};

    % --- 7. 頂点 V プロット ---
    \node[circle, fill=red!85!black, inner sep=1.5pt] at (V) {};
    \node[above right=2pt, red!85!black, font=\footnotesize] at (V) {頂点};

  \end{tikzpicture}
  \caption{放物線 $y = ax^2 + bx$ と直線 $y = -2x + 5$，および囲まれた面積 $S$}
  \label{fig:parabola_and_line}
\end{figure}

与えられた放物線を平方完成すると
\begin{align}
  y = ax^2 + bx = a\left(x+\dfrac{b}{2a}\right)^2 - \dfrac{b^2}{4a} \label{eq:1}
\end{align}
である．この頂点$\left(-\dfrac{b}{2a}, -\dfrac{b^2}{4a}\right)$ が
$y=-2x+5$ 上にあるので，
\begin{align}
  & -\dfrac{b^2}{4a} = \dfrac{b}{a} + 5 \notag \\
  \therefore
  &\quad a = -\dfrac{1}{5} \left( \dfrac{1}{4} b^2 + b \right) \label{eq:2}
\end{align}
である．

題意の面積 $S$ とすると
\begin{align}
  S = \int_{0}^{-\frac{b}{a}} (ax^2 + bx) \,dx
  = -\dfrac{a}{6} \left(-\dfrac{b}{a}\right)^3
  = \dfrac{1}{6} \dfrac{b^3}{a^2}
\end{align}
だから，\cref{eq:2}を代入して
\begin{align}
  S
  &= \dfrac{1}{6} \dfrac{b^3}{\dfrac{1}{25} b^2\left(\dfrac{1}{4}b+1\right)^2} \\
  &= \dfrac{25}{6} \dfrac{b}{\left(\dfrac{1}{4}b+1\right)^2} \\
  &= \dfrac{25}{6} \dfrac{1}{\dfrac{1}{16}b + \dfrac{1}{2} + \dfrac{1}{b}} \\
  &\leqq \dfrac{25}{6} \dfrac{1}{\dfrac{1}{2} + 2\sqrt{\dfrac{1}{16}b \cdot \dfrac{1}{b}}} = \dfrac{25}{6}
\end{align}
である．ただし$b>0$からAM-GM不等式を用いた．
等号成立は $\dfrac{b}{16} = \dfrac{1}{b} \iff b=4$ のときである．

{\bf [解説]}
$b=4$の時 $y=-\dfrac{8}{5}x^2+4x$ となり
\begin{align}
  S = \dfrac{8}{5} \cdot \dfrac{1}{6} \left(\dfrac{5}{2}\right)^3 = \dfrac{4}{3} \dfrac{25}{8} = \dfrac{25}{6}
\end{align}
で解答と一致する．これで検算になる．

\end{document}
